 In \cite{bblr} we introduced a collection of new quiver varieties $\mathcal{M} (\Lambda_n,\vec{v}_c,w_c, \vartheta_c)$, ${n\geq1}$ (see below for the notation);
 for $n\neq 2$ they are not Nakajima's quiver varieties, as the quivers involved are not doubles. We proved that~$\mathcal{M} (\Lambda_1,\vec{v}_c,w_c, \vartheta_c)$ is isomorphic to the Hilbert scheme of points of the total space of $\mathcal{O}_{\mathbb{P}^1}(-1)$, and, in particular, that it is therefore irreducible (as the Hilbert scheme is so~\cite{Fogarty}). For $n\geq 2$ we only proved a weaker result, i.e.,
 that only a certain connected component of $\mathcal{M} (\Lambda_n,\vec{v}_c,w_c, \vartheta_c)$ can be identified with $\operatorname{Hilb}^c(\operatorname{Tot}(\mathcal{O}_{\mathbb{P}^1}(-n)))$. However, as $\mathcal{M} (\Lambda_2,\vec{v}_c,w_c, \vartheta_c)$ is a Nakajima quiver variety, its irreducibility follows from Crawley-Boevey's result, so that one only has to determine whether the varieties $\mathcal{M} (\Lambda_n,\vec{v}_c,w_c, \vartheta_c)$ are irreducible for $n\geq 3$. In this paper we prove this fact, completing the work of \cite{bblr}, actually showing directly
 that $\operatorname{Hilb}^c(\operatorname{Tot}(\mathcal{O}_{\mathbb{P}^1}(-n)))$ is isomorphic to the whole
$\mathcal{M} (\Lambda_n,\vec{v}_c,w_c, \vartheta_c)$. As this technique also works for the case $n=2$ we include it as well.

\section{Some background } \label{background}
The quivers we are going to consider are extracted from the ADHM data for the Hilbert schemes of points of the varieties $\operatorname{Hilb}^c(X_n)$,
 where $X_n$ is the total space of the line bundle $\mathcal{O}_{\mathbb{P}^1}(-n)$,
and, in turn,
the construction of the ADHM data is based on the description of the moduli spaces of framed sheaves on the Hirzebruch surfaces $\Sigma_n$ in terms of monads that was given in~\cite{bbr}.
